\section{Introduction and Background}

Option pricing is one of the central and open questions in the mathematical
finance field. It motivates the advancement of the stochastic process
theory and becomes one of the practical applications of sophisticated
theories. To begin with, options are a type of derivative that gives
the holder the right, but not the obligation, to buy or sell the underlying
asset at a specified price (a.k.a. the strike price) within a specific
period. The name `derivative' implies that those securities are derived
from the origin asset. Although investors can long and short options
for seeking risky profits, one primary feature of options is to obtain
a certain share of assets with a fixed amount of payment in the future.
This helps to control volatility and uncertainty. Thus, the market
maker who sells the option charges the buyer a certain amount of fee
that amounts to the expected risk.

We focus on the European call in this paper. One can only exercise
a European option at the terminal time and a call option gives the
holder the right to buy. The holder usually exercises the European
call if the asset price ends up above the strike price and takes no
action otherwise. If the market is liquid enough, the holder can sell
the asset immediately after obtaining it, making a marginal profit.
Therefore, we define the payoff of a European call as
\begin{equation}
\text{payoff}=\max\left\{ S-K,0\right\} \label{eq:intro-payoff}
\end{equation}
where $S$ is the asset price at the exercise time and $K$ stands
for the strike price. Although it is tempting to price the option
based on this payoff function with consideration of each path probability,
this is usually not the best way since the option is designed to measure
the risk of obtaining the asset. If there is a general trend in the
price process, one can hold a consistent position in the asset and
purchase an option in addition. Thus, the more popular approach is
to manage a portfolio that replicates the option, i.e. the holder
invests the same amount of money as the price of the option in a replication
portfolio and gets the same amount of payoff (Eqn. \ref{eq:intro-payoff})
at the terminal time. The portfolio borrows some extra funds from
an external risk-free source to hold a long position of the underlying
asset and adjusts its position according to the change in the asset
price. This approach is referred to as the replication by portfolio
method.

We conduct a brief review in this section, starting by introducing
the Black-Scholes model which solves the optimal hedging strategy
and the pricing problem at the same time, moving to a literature review
where multiple types of asset pricing are involved, and discussing
how to model transaction costs.

\subsection{Black-Scholes Model}

The Black-Scholes formula (or BS formula or short) is the first formula
that prices the European-style call option. There are two equivalent
ways of calculating the price, either by solving the Black-Scholes-Merton
equation or use the risk-neutral measure. The price of an European
call of time span $T$ and a strike price $K$ is solved by 
\begin{equation}
c_{\text{BS}}\left(t,x\right)=xN\left(d_{+}\left(T-t,x\right)\right)-Ke^{-r\left(T-t\right)}N\left(d_{-}\left(T-t,x\right)\right)\label{eq:bs-price}
\end{equation}
where $0\le t<T$ is the current time and $x$ is the current asset
price, with $r$ being the risk-free interest rate, $\sigma$ being
the volatility, $N\left(y\right):=\frac{1}{\sqrt{2\pi}}\int_{-\infty}^{y}\exp\left(-\frac{z^{2}}{2}\right)\d z$
being the cumulative standard normal distribution, and $d_{\pm}$
defined as
\[
d_{\pm}\left(\tau,x\right):=\frac{1}{\sigma\sqrt{\tau}}\left[\log\frac{x}{K}+\left(r\pm\frac{\sigma^{2}}{2}\right)\tau\right].
\]
The optimal hedging position can be derived as the partial derivative
of Eqn. \ref{eq:bs-price} as 
\begin{equation}
u_{\text{BS}}\left(t,x\right)=N\left(d_{+}\left(T-t,x\right)\right).\label{eq:bs-hedge}
\end{equation}
We leave the details to \cite{shreve2004stochastic,alma991057973644906011}
for more in-depth theory and computation. 

This set of theory surpassed most of the option pricing theories when
it was first published and is still helpful from since. With that
being said, the theory relies on a critical assumption that the agent
is able to continuously rehedge the positions of the replication portfolio.
This is either impractical or costly to most of the investing parties.
In the view of stochastic integration, we use a Riemann-style sum
with left endpoints to approximate the Ito integral. The accumulation
of the temporal discretization errors will affect the final payoff
profoundly.

\subsection{Literature Review}

There have been vast works and literatures on option pricing. A few
exercising strategies for American option are reviewed in \cite{li2009learning}.
The state space is composed of asset price trajectories $\left(S_{0},S_{1},\dots,S_{T}\right)$
and an absorbing state $\boldsymbol{e}$ which serves as the destination
after the option has been exercised. The action space simply contains
two actions: hold and exercise. The only non-zero reward is given
when the option has been exercised at the value specified by the option.
The Q-function is derived as in classical RL problems and it is the
quantity of interest. The first method, LSPI (LS policy iteration),
combines LSTD (LS with TD update) and policy iteration together to
achieve efficient learning. The second and third ones, FQI (fitted
Q-iteration algorithm) and LSMC (Least squares Monte Carlo), take
a DP approach where the exact exercising problem is solved at each
trading moment; the only difference lies in that FDI takes a forward
view in time while LSMC is backward, starting from the exercising
time. In the same spirit, \cite{fathan2021deep} investigates the
efficiency and effectiveness of different parametrizaiton structures.
Three algorithms, namely double deep Q-Learning (DDQN), categorical
distributional RL (C51), and implicit quantile networks (IQN) are
compared where DDQN learns the optimal Q-value function while the
other two algorithms try to learn the full distribution of the discounted
reward. They are tested on empirical data as well as simulated geometric
Brownian motion trajectories.

Different RL architectures can be deployed in this field as well.
The Actor-Critic structure is used in \cite{marzban2021deep} for
Equal Risk Pricing (ERP) in a risk averse setting under the framework
studied in \cite{tamar2015policy}. The concept of $\tau$-expectile
\[
\overline{\rho}\left(X\right)=\arg\min_{q}\tau\bE\left[\left(q-X\right)_{+}^{2}\right]+\left(1-\tau\right)\bE\left[\left(q-X\right)_{-}^{2}\right]
\]
is used to elicit a coherent risk measure. The value function is defined
as the portfolio value under the recursive coherent risk measure realized
by expectiles, i.e.
\[
V_{t}\left(S_{t},Y_{t}\right)=\inf_{\xi_{t}}\overline{\rho}\left(-\xi_{t}^{T}\Delta S_{t+1}+V_{t+1}\left(S_{t+1},Y_{t+1}\right)|S_{t},Y_{t}\right)
\]
with terminal condition $V_{T}\left(S_{T},Y_{T}\right)=F\left(S_{T},Y_{T}\right)$
specified by the option contract. With that being established, we
can apply the policy gradient method where the critic network updates
the value estimate which the actor network can refer to and build
its policy upon. The network uses a classical fully multilayer structure
with alternating activation functions.

A hybrid attempt is carried out in \cite{grassl2010reinforcement}
where the pricing strategy is based on a combination of optimal stopping
and terminal payoff. The idea is that the agent can either hold the
derivative until the terminal time, executing the contract to get
the payoff written, or sell the derivative earlier according to the
price at that particular moment. The reward function is thus simplified
defined as the selling/execution price if such scenario happens. The
value function is parameterized by kernel function approximation and
the algorithm is tested using simulated geometric Brownian paths of
an European call option.

We point out that it is also possible to directly solve the HJB equation
if the associated RL problem is formulated as a control porblem. We
refer to \cite{halperin2021distributional} for more details on distributional
offline continuous-time RL learning algorithms.

\subsection{Trading Cost \label{subsec:Trading-Cost}}

One other important factor in option pricing is the trading cost (or
transaction cost, used interchangeably). The trading cost occurs when
the hedging position of the replication portfolio changes, making
a portion of the invested capital unavailable. Thus, the presence
of trading cost increases the price of the option in general.

A way of understanding and modeling the trading cost is to use the
concept of bid-ask spread. Consider the limit order book of a particular
asset at a given time. Empirical observation shows that the selling
(a.k.a. ask) orders are always above the buying orders (a.k.a. bid)
orders since one wishes to buy low and sell high. Moreover, the orders
are not uniformly distributed: there are fewer orders close to the
mid price (i.e. the arithmetic average of the highest bid and lowest
ask) and more orders away from the mid price. The gap between the
lowest ask and highest bid is called the bid-ask spread. A larger
gap usually implies a higher friction in the market.

In general, the dynamics of the limit order book distribution is a
very challenging problem. However, under the assumption that the replicating
portfolio hedges a small position, we can safely conclude that the
trading cost, defined as the additional cost paid for buying (or selling,
correspondingly) a unit of asset with reference to its mid price,
can be computed as follows
\[
\text{Trading cost (TC)}=S_{\text{ask}}-S_{\text{mid}}=\frac{1}{2}\left(S_{\text{ask}}-S_{\text{bid}}\right)=\frac{1}{2}\text{spread}.
\]
To model the spread, we assume that it scales with the mid price by
a factor of $\epsilon$, characterizing the friction. Thus, the trading
cost to buy/sell $\Delta u$ shares of the asset requires a trading
cost of
\[
\text{TC}\left(\Delta u,S_{\text{mid}}\right)=\frac{\epsilon}{2}S_{\text{mid}}\Delta u.
\]